\section{Failure modes, safety, and governance}\label{sec:failures}


\regime{Control}{how persistence fails, and the governance operations that should prevent it.}
The previous parts established what always-on agents accumulate (substrates, lifecycle operations, continual adaptation) and how the field measures it. The missing counterpart is governance of the state being accumulated. We treat governance and failure modes as two faces of one object. A persistence-induced failure is the observable symptom; a governance mechanism is the lifecycle operation that should have prevented it. Organizing the two together is deliberate. Prior agent-security surveys catalogue threats by attack surface, by component, or by adversary capability \citep{chu2026systematicsurveyagentsecurity,lin2026surveylongtermmemorysecurity}. We instead index every failure to the lifecycle stage at which corrupt state enters and to the invariant (Section~\ref{sec:lifecycle}) it violates, because that connects an attack to the governance operation that should have caught it. A poisoning attack is a write that escaped authority and provenance checks and then persisted. A stale-commitment failure is an update boundary that let an expired authority keep licensing action.

This framing also lets us state the survey's corpus-scoped map precisely. In our $435$-work coded corpus, governance is the least developed region of the lifecycle. By state axis, authority is the rarest at $72$ of $435$ works, recoverability appears in $112$, and provenance in $153$; by lifecycle stage, the return arc thins, with audit at $88$, forget at $66$, and rollback at only $27$. The counts are scoping estimates, but the qualitative skew is large. Fifteen targeted rounds aimed at governance, authority, rollback, formal verification, durable execution, deletion propagation, and the database, distributed-systems, and HCI literatures grew the corpus more than fourfold yet moved the governance fraction only from roughly one-sixth to one-third before late rounds showed diminishing returns under this query frame. We read this as evidence that governance-targeted search did not remove the skew within our frame, not as a proof of field-level saturation. Subarea by subarea, the literature supplies fragments of governance while leaving the lifecycle incomplete. That pattern is the gap: always-on agents are persistent-state systems whose state must be governed over a lifecycle rather than only stored and retrieved.

\subsection{Threat model and governance surface}

Before cataloguing failures, we fix the threat model used here. The attacker need not control the model or the memory store. The typical persistence attack is weaker and harder to notice: an untrusted channel writes state, a trusted channel later retrieves it, and a tool call or shared-memory update gives that state authority it never had. Table~\ref{tab:threat-model} names the principals and channels we assume, what each may write or read, and the recovery question a governed system must answer. The table is intentionally operational rather than cryptographic. It says where the lifecycle must bind authority, provenance, scope, and rollback handles if the later failure taxonomy is to be testable.

\begin{table}[t]
\centering
\scriptsize
\setlength{\tabcolsep}{3pt}\renewcommand{\arraystretch}{1.16}
\begin{tabularx}{\linewidth}{p{0.17\linewidth}p{0.10\linewidth}p{0.24\linewidth}p{0.22\linewidth}Y}
\toprule
Principal or channel & Trusted? & State written or read & Authority surface & Failure to test or recover \\
\midrule
Owner user instruction & Usually & Preferences, permissions, corrections, deletion requests & Direct authority, but only within declared scope and current epoch & Does a later action use the current instruction, not a superseded one? \\
Collaborator or group member & Partial & Shared facts, proposals, edits, delegations & Role-bounded authority over shared state & Can the system surface conflicts and prevent cross-principal scope expansion? \\
Untrusted web, file, email, or tool output & No & Observations, retrieved snippets, tool results, notes embedded in data & No authority to become durable instruction without validation & Can poisoning be quarantined before it writes persistent state or chooses a tool? \\
Memory manager, summarizer, or consolidator & Partial & Summaries, derived facts, indices, embeddings, skill records & Derivative authority inherited from sources & Does consolidation preserve source, scope, and deletion lineage? \\
Agent reflection or self-update loop & Partial & Lessons, skills, policies, acceptance scores & Authority only after validation and rollback registration & Can a bad self-update be traced, de-authorized, and reverted? \\
External API, tool server, or credential broker & Partial & Side effects, credentials, tool schemas, execution receipts & Capability grants and schema-level permission & Can rollback distinguish replay from compensation after an irreversible effect? \\
Administrator or organization policy & Trusted for policy & Retention rules, allowlists, logging, tenant boundaries & Global constraints over user and agent state & Can policy changes narrow authority without silently orphaning old state? \\
\bottomrule
\end{tabularx}
\caption{Threat model for persistent-state failures. A channel may be trusted for one purpose and untrusted for another: an external tool can be trusted to report a weather value but not to write a durable instruction; a collaborator can edit shared state but not expand another user's scope. The recovery column states the check a benchmark or audit should force.}
\label{tab:threat-model}
\end{table}

\subsection{Persistence-induced failure modes}

Episodic agents avoid an entire failure class by resetting. An attack that succeeds within a single task disappears when the task ends; a stale fact cannot govern a future decision because there is no future decision that shares state. Persistence removes this safety net. The failure modes share the same enabling property that makes always-on agents useful: state written in one interaction authorizes behavior in a later one, after the originating context, and often the originating user, is gone. Table~\ref{tab:failures} maps each failure to the lifecycle boundary where corrupt state enters and the invariant it breaks; the surrounding text follows those boundaries in order.

\begin{table}[t]
\centering
\small
\setlength{\tabcolsep}{4pt}\renewcommand{\arraystretch}{1.2}
\begin{tabularx}{\linewidth}{p{0.15\linewidth}YYp{0.12\linewidth}}
\toprule
Boundary & Failure & Invariant violated & Repr.\ work \\
\midrule
Write & Untrusted write poisoning; indirect prompt injection & Authority monotonicity; provenance preservation & \citep{louck2026originbound,joshi2026eywa} \\
Consolidate & Retrieval-key loss; procedure poisoning & Provenance preservation & \citep{ouyang2026memlineage,kumar2026memarchitect} \\
Retrieve and act & Memory distraction; control-flow hijack & Authority monotonicity; scope non-expansion & \citep{xu2026memorysilent,winston2026solveraided} \\
Update & Stale commitment; belief drift & Authority monotonicity & \citep{zhou2026usercorrections,shawn2026pace} \\
Forget and audit & Deletion failure; rollback gap & Deletion propagation; rollback traceability & \citep{zhao2026memorepair,shah2026unlearningmirage} \\
Shared state & Cross-agent propagation & Scope non-expansion & \citep{ren2026gatemem,ravindran2026portablemem} \\
Recovery & Rollback as attack surface & Rollback traceability & \citep{zheng2026acrfence} \\
\bottomrule
\end{tabularx}
\caption{Persistence-induced failure modes, each mapped to the lifecycle boundary where corrupt state enters or activates and to the state invariant (Section~\ref{sec:lifecycle}) it violates. The table treats safety as a property of the whole lifecycle, not a final filter applied after retrieval. The last row is important because recovery, the mechanism meant to repair the others, is itself a fresh boundary at which new corruption can enter.}
\label{tab:failures}
\end{table}

\subsubsection{Write poisoning and injection.}
At the write boundary the core failure is that untrusted observation becomes trusted state. Memory-poisoning and indirect prompt-injection studies show that a harmful record can enter long-term state during an ordinary interaction and activate much later, when the originating context is gone. The threat is not insertion alone but retrieval-aware optimization of what gets inserted so it is reliably surfaced and acted upon \citep{dash2026memorypoisoning}, and the same mechanism scales from single-session prompt manipulation to a cross-session system vulnerability once the payload survives in memories, files, and tool registries after the attacker's session ends \citep{xie2026crosssessioninjection}. Filtering the current prompt cannot undo a write that already landed. The governance response is to push authority and provenance to the write itself. Origin-bound write authority ties a record's right to influence action to a machine-checked, non-malleable account of its source, arguing that content-based and lineage-based trust can be laundered through summarization and so must be cryptographically rather than heuristically bound \citep{louck2026originbound}. Evidence-before-belief designs store immutable source evidence first and derive canonical facts only after validation against typed signals and source support, so a poisoned observation cannot become a belief without leaving an auditable derivation \citep{joshi2026eywa}. These works contribute the validate stage that most memory systems leave implicit: a write is not committed until an authority and a provenance chain license it. What they leave open is coverage. Origin-binding governs the write but assumes the provenance signal is itself trustworthy, and evidence-first storage raises storage and latency cost that no benchmark yet prices. The gap, in thesis terms, is that authority monotonicity and provenance preservation are enforced at the moment of writing by a handful of 2026 systems and assumed away by the rest.

\subsubsection{Retrieval distraction and tool-drift.}
At the retrieve-route-act boundary the failure is that stored state crowds out current intent. Retrieved records can distract the model, override newer evidence, or bias tool selection and ordering. A subtle and well-named form is the silent integration failure: an agent integrates sensitive or irrelevant memory content into its behavior even when no current prompt warrants it, and the integration happens after retrieval, so admitting the right documents does not prevent the wrong influence \citep{xu2026memorysilent}. This sharpens a point the retrieval literature often blurs: retrieval quality is necessary but not sufficient, because the generator can act on retrieved state in ways the retriever never intended. The governance response is to treat the boundary between retrieved state and action as an enforcement point rather than a similarity threshold. Solver-aided execution compiles natural-language tool-use policies into SMT constraints over agent-observable state and tool arguments, then gates each tool call against a precondition check before it executes \citep{winston2026solveraided}; skill-containment analysis bounds, by abstract interpretation and bounded model checking, what a non-deterministic skill can reach in its deterministic execution side \citep{metere2026skillcontainment}. These contribute a runtime authority check at act time, but leave implicit the interaction with memory: a containment proof over a skill says nothing about whether the state that parameterized the skill was authorized to do so. The gap is that authority monotonicity and scope non-expansion are checked, when they are checked at all, over the action's static envelope, not over the dynamic state that selected and parameterized it.

\subsubsection{Stale commitment and belief drift.}
At the update boundary the failure is temporal: old preferences, facts, commitments, permissions, or policies remain binding after they cease to be valid. A preference that was true last month, changed yesterday, should be superseded today, yet a system with no notion of supersession treats the historical value as current. A useful result here is the access-versus-compliance gap: even when a user's prior correction is present and recalled, agents still violate it most of the time, because recall is not compliance \citep{zhou2026usercorrections}. The proposed fix mines corrections into atomic rules compiled into act-time completion gates the agent must pass, again the validate and act stage made explicit rather than assumed. Belief drift is the dual failure during self-update: greedy, score-based acceptance of self-generated updates p-hacks the agent into committing false or harmful state, fixed by an anytime-valid, betting-based commit gate that bounds the false-commit probability \citep{shawn2026pace}. Both works contribute a notion of update as a gated transition rather than an overwrite. Neither closes the representational gap: to mark a value superseded rather than deleted, the store must represent both current and historical truth with authority and timestamps, which many flat-text memories do not. The gap is that authority monotonicity over time, the requirement that only a current, unrevoked authority may license a state-affected action, is rarely representable, so deletion becomes the only available tool for retiring obsolete state.

\subsubsection{Deletion failure and the rollback gap.}
At the forget-and-audit boundary the failure is that deletion is treated as a single operation when it is a cascade. A user asks the system to forget a fact, but that fact may already exist in raw logs, summaries, embeddings, cached prompts, retrieved traces, shared memories, and procedural skills. Deployment-time memorization work makes the residue concrete: derived tiers retain deleted content unless the purge covers the full pipeline \citep{chen2026deploymenttimememorization}, and the recoverability failure is demonstrated directly when supposedly forgotten information resurfaces under multi-hop or aliased queries \citep{shah2026unlearningmirage}. The governance response is cascade-aware deletion. A barrier-first repair contract withdraws stale derived descendants when a source is deleted or invalidated and republishes only validated, predecessor-closed successors, cutting invalidated-memory exposure toward zero \citep{zhao2026memorepair}. This contributes the deletion-propagation invariant in operational form. What stays open is that the contract assumes a tracked derivation graph; systems that consolidate by lossy summarization have already destroyed the lineage the cascade would follow, which is why provenance preservation and deletion propagation are coupled rather than independent.

\subsubsection{Cross-agent propagation.}
Shared state adds a propagation failure absent in single-agent settings. Corrupt state can spread through a shared memory substrate or an inter-agent channel faster than per-agent safety checks can contain it, because each agent trusts the substrate. A governed shared-memory benchmark frames the test directly: can an agent serve legitimate cross-principal requests while enforcing per-role scope and reliably forgetting on deletion request \citep{ren2026gatemem}. Cross-agent memory-interoperability protocols make the substrate explicit, moving persistent memory across heterogeneous vendor models with content-addressable provenance and capability-scoped disclosure \citep{ravindran2026portablemem}, and a workspace-collaboration protocol that delegates a live execution environment rather than only messages widens the shared surface further \citep{nie2026awcp}. The contribution is to name scope as the governing axis: a shared write must narrow or preserve scope, never silently widen it. The gap is that each new shared channel enlarges exactly the authority and provenance surface the propagation failures exploit, and the interoperability protocols assume the receiving model trusts the sender's encodings.

\subsubsection{Rollback as a fresh attack surface.}
The final row of Table~\ref{tab:failures} is the most counterintuitive and among the most relevant to the survey's claim. Recovery, the mechanism meant to repair every failure above, is itself a boundary at which new corruption enters. Because an agent regenerates slightly different requests after a checkpoint restore, retried tool calls become new external requests, enabling duplicate irreversible side effects and the resurrection of already-consumed authority unless irreversible effects are recorded and replayed rather than re-issued \citep{zheng2026acrfence}. This is why rollback traceability is an invariant rather than a convenience feature: a system that can revert its internal store but cannot reason about which external effects were irreversible will, on recovery, re-send a payment or re-grant a revoked permission. The gap is structural. The literature treats rollback as a capability to add, not as a transition that must itself satisfy authority monotonicity and scope non-expansion, and the corpus count makes the consequence stark: with only $27$ of $435$ works exposing any rollback mechanism, the operation most likely to silently re-introduce corruption is the least studied.

\subsection{Authority and permission state}

Authority is the rarest axis in the corpus ($72$ of $435$) and the one whose absence most directly undermines the persistent-state thesis: a record that can influence action without a current, scoped permission is the defining ingredient of poisoning, privilege drift, and cross-user leakage. The governance question is not who owns the memory but which authority currently licenses a given record to affect a given action, and authority must only narrow, never silently expand, as state and delegations accumulate.

One line treats authority as a server-side, monotone policy attribute. Intent-governed authorization makes the user's expressed intent a session-scoped attribute that may only reduce the authority static credentials grant to an agent's tool calls \citep{zhu2026intentgovernedauth}; this is authority monotonicity enforced at the API boundary. A complementary execution-time line attacks the time-of-check-to-time-of-use problem directly: a state-witness profile revalidates preconditions at execution and fails closed when world state has drifted between an action's approval and its delayed execution \citep{zhu2026openport}, which is exactly the stale-authority failure of the previous subsection caught at act time. Delegation is governed compositionally by overlaying recursive-delegation, dynamic-scoping, and resource-attenuation primitives onto existing relational authorization, with formal proofs that delegated authority narrows along the chain \citep{ibrahim2026overlaygovernance}, and extended at the credential layer by scoped agent credentials with accountable delegation chains and natural-language permissions compiled into auditable access control \citep{south2025authenticateddelegation}. A particularly clean formal result imports hardware memory-consistency theory to bound unauthorized operations by execution count rather than time-to-live, so revocation does not depend on agent speed \citep{parakhin2026bureaucracy}.

A second execution-time result makes the monotonicity argument constructive: authority is operationalized as a runtime enforcement mechanism that lets an action execute only when a current authority can be constructed from the system state, with a halt state and recovery loop that suspends execution whenever authority is undefined, so a stale or revoked permission can never license a tool call \citep{fernandez2026authority}. At the coordination layer, lighter-weight declarative manifests cover the open-web case: an agent-permissions.json convention lets site owners declare which agent interactions are allowed, supplying a practical permission-boundary mechanism across heterogeneous agent modalities where no shared credential infrastructure exists \citep{marro2025permissions}. DEMM-Bench makes reconstructability measurable by scoring whether recorded delegation, policy, and tool-firewall decisions are sufficient to rebuild an authority decision and its accountability trace \citep{solozobov2026demm}.

Tool and credential state form a tightly coupled sub-problem, since a stale token or a silently redefined tool is an authority leak by another name. A trusted broker lets an agent invoke API and SSH secrets through schema-checked calls while keeping the raw credential non-exportable, defeating prompt-injection exfiltration \citep{jin2026capseal}; tool definitions are bound to signed, versioned manifests so a persisted grant is invalidated when a tool is silently redefined, the rug-pull case \citep{bhatt2025etdi}; and tool servers are admitted only via offline-signed clearance against a pinned trust root under a deny-by-default allowlist \citep{metere2026attestedtoolserver}. The most thesis-aligned credential result targets the zombie-agent failure: a heartbeat-bound protocol makes a descendant agent's authority self-expire within a provably bounded window once parent liveness ceases, so an agent cannot keep acting on delegated authority after the operator who granted it has shut down \citep{deochake2026heartbeatcred}. This is authority monotonicity given a clock.

This body of work is the governance subarea with the deepest enforcement mechanisms and formal guarantees. What it leaves implicit is the link to memory: these systems govern the credential and the tool call, but say little about the persistent state that decides which credential to use or which tool to call. A monotone-authority API does not help if poisoned memory chooses to invoke an authorized but contextually wrong tool. The gap, then, is that authority is governed at the action boundary by a small, sophisticated literature and almost never co-governed with the memory whose retrieval triggers the action. Revocation dynamics, especially partial-delegation revocation and revocation cascades, remain largely unaddressed.

\subsection{Provenance, audit, and lineage}

Provenance is better covered than authority ($153$ of $435$) but is dominated by static logging rather than the derivation-aware lineage that deletion and rollback actually require. The governing question is whether a system can answer, for any stored fact or any action, where it came from, how it was transformed, and which other state descends from it.

Three mechanistic families recur. Append-only ledger designs treat memory as an immutable, signed history: a Merkle-indexed ledger with privilege-lattice access control makes revisions protocol-blocked and historied rather than destructive \citep{wright2025immutablemem}, and a signed delegation-provenance chain binds each agent action to an originating human authorization and scope for offline audit \citep{dalugoda2026hdp}. The append-only pattern has a deeper lineage in the systems literature that the agent work is rediscovering: blockchain-backed audit logging supplies an established design for tamper-proof, cryptographically secured action records whose state transitions cannot be retroactively altered \citep{ahmad2018blockaudit}, and smart-contract-backed audit trails extend this to autonomous agents by atomically recording each action, prediction, and explanation to a ledger so decision provenance is immutable and regulator-auditable \citep{wang2026auditor}. A broader honest-computing framework generalizes the requirement, establishing provenance through cryptographically backed data-lineage tracking and auditable custody chains that apply directly to persistent agent memory states and their transformations \citep{guitton2024honest}. Derivation-graph designs go further by tracking history and descent: MemLineage attaches a weighted derivation DAG and Merkle log and enforces an untrusted-path-persistence invariant that refuses actions descending from external ancestors \citep{ouyang2026memlineage}, while portable-memory and robot-policy-lineage designs carry content-addressable provenance across vendor models and across policy-iteration cycles, respectively \citep{ravindran2026portablemem,luo2026robolineage}. The same derivation discipline is what self-improving agents need to make learned behavior auditable: recent work reframes capability updates as traceable artifacts by recording verifiable provenance of skill acquisition through structured skill graphs and cryptographically backed audit logs, rather than as opaque parameter shifts \citep{huang2025audited}. Formal-monitoring designs make provenance checkable at runtime: Causal Past Logic gives a past-time temporal logic with a provably correct vector-clock monitor for guards over causally-visible stored state in distributed agent workflows \citep{bollig2026causalpast}. A complementary diagnostic line shows why provenance must be structured, not flat: when flat-text memory loses source attribution, agents hallucinate authority and misattribute facts, a failure named provenance-role collapse \citep{jin2026provenancerolecollapse}. That the discipline is far from settled even in evaluation is itself documented: a systematic audit of agent-benchmark disclosures finds reproducible provenance of evaluations still nascent, with mean audit scores of $0.38$ out of $1.0$ on core dimensions such as harness specification and inference-cost reporting \citep{moghadasi2026audit}.

The tension within this subarea is itself a finding. Immutability and deletion are in direct conflict: an append-only ledger that makes revisions protocol-blocked cannot, by construction, honor a right-to-be-forgotten request without breaking its own invariant \citep{wright2025immutablemem}. Derivation graphs resolve part of this by enabling cascade deletion over lineage, but assume an acyclic, fully-tracked graph and do not handle benign cycles introduced by cross-agent delegation \citep{ouyang2026memlineage}. The gap is that provenance is mostly logged for accountability, not maintained as the substrate that deletion propagation and rollback traceability depend on; few works treat lineage as a live, queryable, prunable structure rather than a write-once record, and provenance preservation under lossy consolidation, where the lineage is destroyed exactly when it is most needed, is the single thinnest cell.

\subsection{Privacy, consent, and scope}

Privacy is where governance work is densest in absolute terms, yet it remains skewed toward demonstrating leakage rather than enforcing scope. The persistent-state framing reclassifies privacy from a training-time property to a deployment-time, cross-session one: the risk is that an always-on agent silently writes, retains, and later exposes user state it was never explicitly authorized to keep, rather than only that a model memorized its training data \citep{carlini2021extractingtrainingdata}.

The leakage-demonstration literature is now precise about the attack surface. Membership-inference frameworks target persistent chat-agent memory units directly \citep{chen2026mrmmia}, an interrogation attack infers datastore membership with roughly thirty natural-language queries that evade query-rewriting defenses \citep{naseh2025riddle}, mobile-agent memory leaks private records under extraction \citep{wang2025mobileagente}, and deployment-time memorization studies show how deployed agents absorb sensitive interaction data \citep{chen2026deploymenttimememorization}. A particularly consequential audit of $2{,}050$ real persistent-memory entries finds that $96\%$ are silently system-created rather than user-authorized, motivating an attribution shield for transparency over autonomously written user state \citep{dash2026selfportrait}. The enforcement-side literature is thinner. Differential privacy for memory spends budget only on sensitive tokens to bound leakage from an external store \citep{koga2024dprag}; a watermark hidden in latent memory-write decisions lets an owner prove provenance of a leaked snapshot at near-zero utility cost \citep{zhang2026memmark}; and a secure exchange protocol gives cross-session sharing under AES-256-GCM access control \citep{masoor2025samep}. Two HCI studies supply the consent dimension the mechanism work omits: interviews reveal that users hold incomplete mental models of RAG-based memory and want granular review, edit, delete, and categorize control \citep{zhang2025ragmemoryperceptions}, and an analysis of $22$ agentic systems derives UI patterns with regulatory potential, positioning the interface as a governance surface that should make agent memory user-editable \citep{feng2025regulatoryui}.

Contextual integrity gives this subsection a positive criterion for scope: private state must not leak, and information flows must remain appropriate to the social and task context in which the state was written. CIMemories turns this criterion into a benchmark for persistent memory, asking whether memory-augmented assistants preserve contextual integrity across composed situations \citep{mireshghallah2025cimemories}. Contextualized privacy defense for LLM agents makes a complementary mechanism claim, treating privacy intervention as a context-dependent runtime decision rather than a static filter \citep{wen2026contextualizedprivacy}. In our terms, both lines are scope-non-expansion work: a record written in one context must not silently authorize use in another.

The data-centric privacy survey makes the gap explicit. It organizes risks by the data surfaces an agent touches and finds that information-flow control is the only governance covering cross-session inference leakage, while no benchmark exercises an agent across all its data surfaces under a single declared privacy policy \citep{lahjouji2026privacysurvey}. This is the scope gap in operational form. Scope non-expansion requires that state written under one user, task, or consent context not silently act under another, yet the corpus measures leakage after the fact far more than it enforces scope before it. The gap, in thesis terms, is that privacy is studied as an attack to demonstrate rather than a scope invariant to maintain jointly across the write, organize, retrieve, and share stages, and consent is largely confined to the interface layer rather than bound to the stored record.

\subsection{Public production memory controls}

Public documentation for deployed assistants gives a useful deployment-facing cross-check on the survey's abstractions. These systems are not research benchmarks, and the available documentation is not enough to score their internal enforcement. It does show which controls providers expose to users, administrators, and projects. Table~\ref{tab:production-controls} reads those controls through the same axes used in the corpus.

\begin{table}[t]
\centering
\scriptsize
\setlength{\tabcolsep}{3pt}\renewcommand{\arraystretch}{1.16}
\begin{tabularx}{\linewidth}{p{0.17\linewidth}p{0.24\linewidth}p{0.25\linewidth}Y}
\toprule
System & Persistent state exposed in docs & User or admin controls & Persistent-state reading \\
\midrule
ChatGPT Memory \citep{openai2026chatgptmemoryfaq} & Saved memories and reference-chat-history signals that can personalize later responses & Users can view, update, delete saved memories, disable memory, use temporary chats, and delete chats, but deletion across memory, chat history, files, and connected apps is operationally separate & Strong user-facing edit/delete affordance; weaker public account of derived-state provenance and deletion propagation across all downstream uses \\
Microsoft 365 Copilot Memory \citep{microsoft2026copilotmemorysupport} & Work-context memories and user preferences used for Microsoft 365 Copilot personalization & Users can ask Copilot what it remembers, update or remove memories, and disable memory; organizational controls mediate availability & Treats enterprise memory as governed tenant state, but the public surface emphasizes user control more than rollback or recovery of actions conditioned on old memories \\
Claude Code memory \citep{anthropic2026claudecodememory} & Markdown memory files at project, user, and managed-policy scopes, loaded into the coding context & Users and organizations edit memory files directly; hooks and permissions provide separate enforcement points & Makes scope visible through file location and hierarchy, but memory is instruction context unless paired with tool gates and audit hooks \\
Gemini Enterprise personalization and memory \citep{googlecloud2026geminienterprisememory} & Profile, conversation history, connected sources, and saved memories for enterprise personalization & Administrators configure personalization and memory; users can manage saved memories depending on organizational settings & Exposes the split between user memory and enterprise data connectors, a concrete instance of scope and authority separation across data surfaces \\
\bottomrule
\end{tabularx}
\caption{Publicly documented production controls mapped to the survey axes. The table does not claim to evaluate internal product behavior. It uses official documentation to show that deployed assistants already expose persistent-state controls, while rollback traceability, provenance of derived memories, and cross-surface deletion remain less visible in public interfaces.}
\label{tab:production-controls}
\end{table}

The deployment lesson is modest but important. Production systems already treat memory as governable state: users can inspect or delete some memories, organizations can bound what memory is available, and project-scoped memory can be separated from user-scoped memory. At the same time, the public control surface usually stops at editing or disabling memory. It rarely exposes the lineage of derived memories, the permission epoch under which a memory was written, or the rollback handle for an external action that memory influenced. That is the distinction this survey makes between memory control and persistent-state governance.

\subsection{Deletion propagation and the right to be forgotten}

Deletion deserves its own treatment because it is the operation where the persistent-state thesis bites hardest: in an episodic system deletion is free, since nothing persists, whereas in an always-on system a single forget request must propagate across every derived tier or it is only a retrieval edit. The recoverability axis ($112$ of $435$) is moderately covered, but almost entirely on the wrong side of the operation, as machine unlearning over model parameters rather than cascade deletion over agent memory.

The prerequisite problem, identifying what to forget, is addressed by GDPR and right-to-be-forgotten work that supplies a benchmark for the deletion target most mechanism papers assume away \citep{staufer2025whatshouldllmsforget}. The mechanism literature then splits by substrate. Parameter-resident approaches are surveyed comprehensively for removing data influence from weights \citep{nguyen2024machineunlearningsurvey}, but, as the survey itself notes, target weights rather than the retrieval indices and external stores agent memory actually uses. Within that line, recent methods sharpen what efficient parameter-level erasure requires: one frames unlearning as an information-theoretic boundary adjustment that minimizes the gradient of influence functions around deleted datapoints, making removal tractable without full retraining \citep{foster2024information}, while a complementary result argues that for LLM-based agents deletion must operate at the semantic embedding level rather than on token outputs, since otherwise a prompt-rephrasing attack recovers the supposedly unlearned concept from internal representations \citep{spohn2025align}. A theoretical barrier circumscribes the whole enterprise: in multi-stage training pipelines, unlearning is path-dependent, so deletion propagation that achieves erasure equivalent to a retrained agent must track the full training history and stage order, which most always-on stores do not record \citep{yu2025impossibility}. Retrieval-resident approaches operationalize the forget operation on RAG memory by removing documents and evaluating deletion completeness \citep{wang2024unlearningmeetsrag}, and a membership-leakage study establishes the failure mode for a persistent retrieval store directly \citep{anderson2024ismydata}. A database-centric account reframes the requirement for persistent agent stores: deletion must address logical and physical removal as well as inference leakage through derived data and through the deletion pattern itself, which is why semantic deletion guarantees that bound post-erasure recoverability are the right target rather than mere record removal \citep{chakraborty2026deletion}. The verification dimension is surveyed for the first time as a structured problem, mapping methods that confirm deletion actually succeeded \citep{xue2025unlearningverificationsurvey}, which is the recoverability axis in evaluative form. Formal verification supplies one principled route: framing unlearning as a probabilistic hypothesis test over backdoor-instrumented models certifies erasure with high confidence and so could give right-to-erasure compliance an auditable proof rather than an assertion \citep{sommer2020probabilistic}. That verification is itself adversarially fragile, however: a malicious provider can covertly retain unlearned data while still passing the verification check, so a deletion-propagation guarantee built only on after-the-fact verification can be silently circumvented \citep{zhang2024fragile}. The strongest cascade result, the barrier-first repair contract that withdraws derived descendants and republishes only validated predecessor-closed successors, comes from the state-governance line discussed above \citep{zhao2026memorepair}, and a policy layer that filters stale zombie memories during decay supplies the routine, non-adversarial face of the same need \citep{kumar2026memarchitect}.

The collision with provenance is the defining structural problem. An immutable ledger maximizes auditability but makes deletion protocol-impossible \citep{wright2025immutablemem}; a derivation graph enables cascade deletion but presumes lineage that lossy consolidation destroys \citep{ouyang2026memlineage,zhao2026memorepair}. The gap is that no corpus work formalizes delete propagation across dependent entries, indices, summaries, and downstream actions with cryptographic proof of absence, and verification of deletion in retrieval-based agent memory, as opposed to parametric models, remains largely unaddressed. Deletion propagation appears here as an invariant because the literature usually treats deletion as a feature of one tier rather than a property of the whole state.

\subsection{Rollback and recovery}

Rollback is the most under-served capability in the coded corpus and one of the survey's main open problems. Only $27$ of $435$ works expose any rollback mechanism, the rollback share was $0\%$ before 2025 and reached only $9.5\%$ in 2026, and we found no coded work that reports recovery success or cost after corruption. Recovery closes the lifecycle loop: every other governance mechanism detects or prevents corrupt state, but rollback is what repairs the state-affected decisions a now-discovered bad record already caused. Its near-absence means the works we coded can increasingly notice that state went wrong but much less often undo the consequences.

The works that do address recovery cluster into two thin lines. Transactional designs import database guarantees: a Saga-style model gives multi-agent planning workflow-wide consistency through compensable execution, automated compensation, independent validation agents, and modular checkpointing under relaxed-ACID semantics \citep{chang2025sagallm}, bridging the classical distributed-systems saga pattern with its compensating transactions to agentic planning, and a graph-native versioned memory proves its update and contraction operations satisfy the AGM belief-revision postulates, importing classical knowledge-representation rationality guarantees into cross-session persistent state \citep{park2026kumiho}. A complementary task-level design supplies the missing transaction boundary for tool use: semantic transactions wrap multi-step agent tool invocations in commit-rollback-recovery scopes with tracked effect lineage and reversible local state, so a partially completed task can be cleanly contracted rather than left in a torn state \citep{chen2026cordon}. These contribute principled undo and consistent contraction. Their limitation is the assumption that compensation handlers exist and that original state can be reconstructed, which fails when an action had an irreversible external effect. The second line confronts that failure head-on: the recovery boundary is itself an attack surface, because regenerated requests after a checkpoint restore become new external requests that duplicate irreversible effects or resurrect already-consumed authority unless those effects are recorded and replayed rather than re-issued \citep{zheng2026acrfence}. This is the strongest recovery result in the corpus because it shows rollback is not a safe escape hatch but a transition that must itself satisfy authority monotonicity and scope non-expansion.

A third, more infrastructural line is now emerging that the agent literature can borrow from rather than reinvent. The durable-execution side supplies the substrate beneath any agent-specific protocol: device-resident checkpointing recovers GPU state and restarts long-running inference without reconstructing KV caches \citep{gan2026concordia}, write-ahead logging with quorum-replicated durability provides foundational low-latency checkpoint patterns decoupled from agent semantics \citep{kuschewski2026btrlog}, and checkpoint-resume orchestration lets long-running multi-agent systems pause on upstream failure and resume converged work across process restarts with reduced recovery overhead \citep{annapureddy2026prima}. These answer the recovery-time question by making fast, durable restart cheap, though none yet reasons about which restored effects were irreversible. The measurement side is finally being instrumented as well, addressing the corpus-wide absence of recovery metrics noted above: MEMPROBE establishes that successful task completion and recoverable memory state are distinct capabilities, turning recovery accuracy into a quantifiable governance metric rather than an assumed property \citep{ma2026memprobe}, and a mean-time-to-recovery metric adapted for agent coherence (MTTR-A) extends classical dependability measures to distributed multi-agent systems by quantifying the time to detect and recover from reasoning drift \citep{or2025mttralm}. Together these begin to supply the recovery-time and recovery-point objectives whose total absence we flag next.

The gaps here compound. There is no quantification of recovery-time or recovery-point objectives, no governance-overhead or latency budget, and no longitudinal production case study; transactional consistency is almost entirely single-agent, leaving distributed and multi-agent rollback semantics absent; and the open instrument named in the research agenda is authority-scoped rollback of tool-originated state once the authority that licensed the original action has lapsed, a case we found no current benchmark scoring. Stated against the thesis, rollback traceability is the invariant most directly tied to repair: every action should carry a handle back to the records that justified it so that a later-discovered bad record yields a bounded, repairable set of decisions. In our corpus this is enforced by a handful of 2026 systems and rarely measured.

\subsection{Synthesis: governance as the missing half of the lifecycle}

Read together, the six subareas tell one story. The works we coded contain sophisticated, sometimes formally verified machinery for the parts of governance closest to the action boundary, monotone authorization, capability containment, attested tool servers, and a substantial descriptive literature on the attacks that persistence enables. Across these enforcement results, validation works best as a live runtime gate rather than a static policy check; recent work makes this explicit. Behavioral contracts integrate preconditions, invariants, governance policies, and recovery as executable specifications enforced at runtime to bound behavioral drift in autonomous agents \citep{bhardwaj2026contracts}, supplying a frame in which our validate stage and every invariant become checkable obligations rather than aspirations. Concrete instances populate that frame: a dedicated runtime authority-control layer gates tool invocations and tracks source and target trust to prevent authority confusion independent of prompt-only defenses \citep{qin2026airguard}; the halt-and-recover authority construction discussed above guarantees no action executes without constructible authority \citep{fernandez2026authority}; and capability mutation at runtime is bounded by explicit validation, traceability, and rollback constraints so that adaptation stays observable rather than unconstrained \citep{garralda2026governed}. Two framing results bound when such gates should be used: a governance rubric reserves runtime guardrails for controls that are observable, determinate, and time-sensitive, assigning the rest across architecture, policy, escalation, and audit \citep{koch2026governance}, and a proactive monitor learns behavioral dynamics to predict unsafe trajectories and intervene before a violation manifests rather than after \citep{wang2025probguard}, which complements the reactive, post-hoc detection that dominates the corpus. What remains thin in our coded corpus is the return-arc machinery that closes the lifecycle: cascade-aware deletion that survives lossy consolidation, lineage maintained as a live structure rather than a write-once log, scope enforced across all of an agent's data surfaces under one policy, and rollback that is itself governed rather than a fresh attack surface. The numbers make the asymmetry concrete as a scoping estimate: authority at $72$, recoverability at $112$, audit at $88$, forget at $66$, and rollback at $27$, all out of $435$, with late targeted search rounds showing diminishing returns under our query frame. Here the skew is a safety concern as well as a coverage statistic: every failure mode in Table~\ref{tab:failures}, from a poisoned write that activates after its author is gone to a recovery that re-sends a payment, is enabled by a missing return-arc operation. The five invariants, authority monotonicity, scope non-expansion, deletion propagation, provenance preservation, and rollback traceability, are therefore not abstract desiderata but the named properties whose absence each documented attack exploits, and which the existing literature names, occasionally enforces, and rarely measures end to end.
